\section{Metrics and cost model}
\label{research-design-and-methodology:metrics-and-cost-model}

The experiment reports each workload and configuration along three outcome dimensions: wall-clock time, network transfer, and operational cost. These correspond to the read-performance, transfer, and cost terms named in the research questions, so that every comparison resolves into the same axes regardless of the configuration under test. Network transfer is treated as a primary outcome rather than an incidental statistic, closing the accounting gap set out in Chapter~\ref{ch:related-work}: prior spatial benchmarks seldom measure the bytes a query moves, yet a cloud benchmark is incomplete unless it reports end-to-end performance together with the pricing that performance drives \parencite{folkerts2012_benchmarking_in_the_cloud}. The remainder of this section defines each dimension in turn and develops the cost model that converts the first two into the third.

\subsection{Primary metrics}
\label{research-design-and-methodology:metrics-and-cost-model:primary-metrics}

Wall-clock time is the headline metric and measures the elapsed interval from the moment a query is submitted to the moment its result is available, as defined for the distributed lifecycle in Section~\ref{background:distributed-spatial-processing:the-lifecycle-of-a-distributed-spatial-query}. For the single-node configurations the measurement boundary is the client process that issues the query, so the figure spans request construction, the engine's own execution, and the return of the result to the client. For the distributed configuration the boundary is the same client, but the elapsed interval now contains the full cluster-side lifecycle, from driver-side planning through parallel reads and shuffles to final collection. Reporting one elapsed time at a common boundary keeps the configurations comparable even though what happens inside that interval differs sharply between them.

Network transfer is measured as two separate quantities, bytes received and bytes sent, recorded at the client or host boundary where the workload issues its reads and collects its results. The two directions are kept apart because a configuration that pushes computation to the data moves different inbound and outbound volumes than one that pulls raw features back to the client, and the split is itself a signal about where work is performed. Operational cost is the third primary dimension. It is defined here only as the monetary charge attributable to running a workload on metered cloud resources, and its decomposition into resource terms is developed in Section~\ref{research-design-and-methodology:metrics-and-cost-model:cost-model}.

\subsection{Auxiliary metrics}
\label{research-design-and-methodology:metrics-and-cost-model:auxiliary-metrics}

Two further quantities are recorded for every run but kept in a diagnostic role rather than reported as outcomes. \acrshort{cpu} time, split into user and system components, indicates how much of an elapsed interval was spent computing rather than waiting on storage or the network, which separates a compute-bound configuration from a transfer-bound one. Result cardinality, the number of rows or features a query returns, records how much data each query actually produced, so that a difference in time or bytes can be checked against the size of the result rather than mistaken for an effect of the format or engine. Neither quantity ranks the configurations; both exist to explain the primary results where an elapsed time or a byte count would otherwise be ambiguous.

\subsection{Distributed phase metrics}
\label{research-design-and-methodology:metrics-and-cost-model:distributed-phase-metrics}

The distributed configuration admits a finer measurement than a single elapsed time, because its wall-clock interval is the sum of distinct phases that scale differently with the workload. The benchmark therefore decomposes distributed wall-clock time into the three phases established in Section~\ref{background:distributed-spatial-processing:the-lifecycle-of-a-distributed-spatial-query}: the executor read time spent fetching input partitions from object storage, the shuffle volume moved between executors when an operation crosses a stage boundary, and the driver collection time spent returning the result set at the end of the query. Attributing elapsed time and transferred bytes to these phases is what lets an observed difference against a single-node engine be traced to a specific mechanism, such as a costly shuffle or a large collection, rather than to the bare fact of distribution.

Each phase term corresponds to a concrete quantity exposed by the Spark runtime: executor read time and shuffle volume are read from the metrics the runtime emits for input and shuffle stages, and driver collection time is taken from the action that returns results to the driver. This section fixes the conceptual correspondence between the thesis-level phases and the runtime quantities that stand in for them; the instrumentation that captures those quantities, including how the runtime's metric stream is intercepted and recorded, is an implementation concern deferred to Chapter~\ref{ch:system-architecture-and-implementation}.

\subsection{Cost model}
\label{research-design-and-methodology:metrics-and-cost-model:cost-model}

Operational cost follows from the cloud cost model of Section~\ref{background:cloud-native-data-architecture:cloud-cost-model}, in which storage, compute, operations, and bandwidth are each metered and priced along a separate dimension. The model expresses the cost of a workload as the sum of four resource terms: compute, charged for allocated processing capacity over the time it is held; storage, charged per gigabyte-month and prorated to the fraction of a month the data occupies; operations, charged per $10{,}000$ transactional requests against object storage; and network egress, charged per gigabyte of outbound traffic. Compute is accounted at per-second granularity rather than the coarser per-minute increment used for illustration in Section~\ref{background:cloud-native-data-architecture:cloud-cost-model}, because the short-lived configurations measured here would otherwise have their cost rounded well above what they consume.

Each configuration is charged only for the resource terms it actually consumes, so the cost expression differs by configuration: a single-node engine accrues compute for the client instance and, where its data sits in a managed relational store, the charge for that store; the distributed engine accrues compute for its cluster across the cluster's lifetime; and any configuration reading from object storage accrues storage, operation, and egress terms. The egress term is treated as near-zero under one modeling assumption, that transfer remaining within a single region and tenant is not billed as outbound traffic, so a charge arises only when the distributed engine moves data across a region boundary. The concrete services that realize each term, the pinned 2026 unit prices, and the offline computation that combines them into a per-workload figure are implementation details, specified in Section~\ref{system-architecture-and-implementation:measurement-instrumentation-and-cost-model:cost-model-implementation}.